\section*{Original Section 4: Lower Bound (journal pp.41--44)}
We prove Theorem 1.10 in this section. The main tool is the notion of Chow varieties, which parameterize projective subvarieties. More precisely, they parametrize a generalization of projective subvarieties, called (effective) algebraic cycles on a projective space.
\paragraph{Algebraic cycles.}
An algebraic $r$-cycle (or simply $r$-cycle) on $\PP^n$ is a formal linear combination $\mathcal D=\sum c_i\mathcal V_i$ of finitely many irreducible subvarieties $\mathcal V_i\subseteq\PP^n$ of dimension $r$, where the coefficients $c_i$ are integers. The degree of $\mathcal D$ is $\deg(\mathcal D):=\sum c_i\deg(\mathcal V_i)$. The support of $\mathcal D$ is $\supp(\mathcal D):=\bigcup_{c_i\ne0}\mathcal V_i$. An $r$-cycle is effective if all its coefficients are nonnegative. Denote by $C(r,d,n)$ the set of all effective $r$-cycles of degree $d$ on $\PP^n$.
\paragraph{Chow varieties.}
Let $k\in\{0,1,\ldots,n-1\}$ and $r=n-k-1$. The definition of Chow forms naturally extends to effective $r$-cycles. Namely, for an effective $r$-cycle $\mathcal D=\sum_{i=1}^r c_i\mathcal V_i$ of degree $d$ on $\PP^n$, define the Chow form of $\mathcal D$ to be $R_{\mathcal D}:=\prod_{i=1}^r R_{\mathcal V_i}^{c_i}$.
Note that $R_{\mathcal D}$ is a vector in $(R/I)_d$ and is uniquely determined by $\mathcal D$ up to scalars. Write $[R_{\mathcal D}]$ for the point in $\PP(R/I)_d$ represented by $R_{\mathcal D}$. Then we have map $\psi:C(r,d,n)\to\PP(R/I)_d$, given by
\[\psi:\mathcal D\mapsto[R_{\mathcal D}],\]
called the Chow embedding of $C(r,d,n)$. Indeed, it embeds $C(r,d,n)$ in $\PP(R/I)_d$ as a projective subvariety, as stated by the following theorem.
\begin{result}{Theorem 4.1 (Chow \& van der Waerden 1937)}
The map $\psi$ is injective and its image is Zariski-closed.
\end{result}
A proof can also be found in Gelfand et al. (1994, Chapter 4). We identify $C(r,d,n)$ with its image under $\psi$ and view it as a projective variety. This variety is called the Chow variety of effective $r$-cycles of degree $d$ on $\PP^n$.
\begin{result}{Example 4.2}
Let $V$ be the subspace of homogeneous polynomials in $\FF[X_1,\ldots,X_{n+1}]$ of degree $d$. Then $C(n-1,d,n)$ is simply the projective space $\PP V$ (see Example 3.9).
\end{result}
\begin{result}{Example 4.3}
$C(r,1,n)$ is the Grassmannian $\GG(r+1,n+1)$ (or $\GG(r,n)$) embedded in $\PP^{\binom{n+1}{r+1}-1}=\PP^{\binom{n+1}{k+1}-1}$ via $\phi^\vee$ (see Example 3.10).
\end{result}
\paragraph{The dimension of Chow varieties.}
When $d=1$, the Chow variety $C(r,d,n)$ is just the Grassmannian $\GG(r+1,n+1)$ (see Example 4.3) and its dimension is well known to be $(r+1)(n-r)$ (Harris 1992). When $d>1$, the dimension of $C(r,d,n)$ was determined by Azcue (1992) and independently by Lehmann (2017). We state their result as follows.
\begin{result}{Theorem 4.4 (Azcue 1992; Lehmann 2017)}
For $d>1$ and $0\le r<n$, the dimension of $C(r,d,n)$ is
\[\max\left\{d(r+1)(n-r),\binom{d+r+1}{r+1}-1+(r+2)(n-r-1)\right\}.\]
\end{result}
This theorem was previously proved by Eisenbud \& Harris (1992) for the special case $r=1$.
\begin{result}{Remark 4.5}
To prove Theorem 1.10, we only need a lower bound for the dimension of the Chow variety, which is much easier to prove than Theorem 4.4. Indeed, it is not difficult to see that $d(r+1)(n-r)$ is the dimension of the space of unions of $d$ $r$-subspaces of $\PP^n$, and $\binom{d+r+1}{r+1}-1+(r+2)(n-r-1)$ is the dimension of the space of degree-$d$ hypersurfaces in $(r+1)$-subspaces of $\PP^n$.
\end{result}
\paragraph{Lower bound via dimension counting.}
We now restate Theorem 1.10 and prove it using a dimension counting argument.
\begin{proof}[Proof of Theorem 1.10]
Consider an arbitrary $k$-subspace $W\in\mathcal H$. We may think of each point in $\PP(R/I)_d$ as a homogeneous polynomial of degree $d$ in Pl\"ucker coordinates modulo scalars and the ideal $I$ of Pl\"ucker relations. We know Pl\"ucker coordinates always satisfy the Pl\"ucker relations. So it makes sense to talk about if a point in $\PP(R/I)_d$ vanishes at $\phi(W)$ or not, as it does not depend on the choice of the homogeneous polynomial representing this point. Note that the constraint of $p\in\PP(R/I)_d$ vanishing at $\phi(W)$ is a linear equation in the homogeneous coordinates of $p$. So the set of points in $\PP(R/I)_d$ vanishing at $\phi(W)$ is a hyperplane of $\PP(R/I)_d$, which we denote by $H_W$.
Let $r=n-k-1$. Assume $|\mathcal H|\le\dim(C(r,d,n))$. Then we have
\[C(r,d,n)\cap\bigcap_{W\in\mathcal H}H_W\ne\emptyset\]
since taking the intersection with a hyperplane reduces the dimension of a projective subvariety by at most one (Lemma 2.7 or Lemma 2.9). So there exists an effective $r$-cycle $\mathcal D=\sum_{i=1}^s c_i\mathcal V_i\in C(r,d,n)$, where $c_1,\ldots,c_s>0$, such that $\psi(\mathcal D)=[R_{\mathcal D}]$ vanishes at $\phi(W)$ for all $W\in\mathcal H$.
Let $\mathcal V=\supp(\mathcal D)=\bigcup_{i=1}^s\mathcal V_i$. Note $\mathcal V\in\mathcal F$ since $\deg(\mathcal V)\le\deg(\mathcal D)=d$. For all $W\in\mathcal H$, we know $R_{\mathcal D}=\prod_{i=1}^sR_{\mathcal V_i}^{c_i}$ vanishes at $\phi(W)$, or equivalently, $R_{\mathcal V}=\prod_{i=1}^sR_{\mathcal V_i}$ vanishes at $\phi(W)$. This implies $\mathcal V\cap W\ne\emptyset$ for all $W\in\mathcal H$. As $\mathcal V\in\mathcal F$, this contradicts our assumption about $\mathcal H$. We conclude
\[|\mathcal H|\ge\dim(C(r,d,n))+1.\]
The dimension of $C(r,d,n)$ is $(r+1)(n-r)$ when $d=1$ and is given by Theorem 4.4 when $d>1$. Plugging in $r=n-k-1$ proves the theorem.
\end{proof}
